% GENERATED FILE. Edit canonical lessons, then run npm run generate:books.
% canonical-source-hash: fnv1a64:01c5d281b298fb8a
% canonical-lessons: TA-C104-language, TA-C104-english, TA-C104-wedding, TA-C104-address, TA-C104-birthday

\chapter{Language, Wedding, Birthday}
\label{ch:ta-language-wedding-birthday}
\hlchaptermodality{104}

\begin{chapteropening}
\textbf{By the end of this chapter:} \emph{I can talk about languages and English, a wedding, an address and a birthday.}

Complete the last lesson of chapter 104: i can talk about languages and English, a wedding, an address and a birthday.
\end{chapteropening}

\section[\texorpdfstring{\ta{மொழி} (moḻi)}{மொழி (moḻi)}]{\ta{மொழி} (moḻi) — a language}

\label{lesson:TA-C104-language}



\begin{quote}
Before the new one: say the Tamil for a family name, then the Tamil for a pet name.
\end{quote}

\subsection*{You'll want to know: \ta{மொழி}}
\textbf{\ta{மொழி}} (\emph{moḻi}) — \textquotedblleft{}a language\textquotedblright{}.

\textbf{\ta{தமிழ்} \ta{என்} \ta{மொழி}} — \textquotedblleft{}Tamil is my language\textquotedblright{}. \textbf{\ta{தாய்மொழி}} is \textquotedblleft{}mother tongue\textquotedblright{}.

\textbf{In the family, and next door}

\noindent
\begin{tabularx}{\linewidth}{@{}>{\raggedright\arraybackslash}X>{\raggedright\arraybackslash}X@{}}
\toprule
\textbf{} & \textbf{word} \\
\midrule
Kannada & \ka{ಭಾಷೆ} (\emph{bhāṣe}) \\
Telugu & \te{భాష} (\emph{bhāṣa}) \\
Malayalam & \ml{ഭാഷ} (\emph{bhāṣa}) \\
Hindi (neighbour) & \dv{भाषा} (\emph{bhāshā}) \\
English & a language \\
\bottomrule
\end{tabularx}

\begin{grammarlens}[title={What you've built}]
A language.
\end{grammarlens}

\subsection*{Your turn}
\begin{itemize}
\raggedright
  \item \emph{Say it:} \emph{moḻi}
  \item \emph{Say it:} \emph{moḻi}, once more
  \item \emph{Say it:} say \emph{tamiḻ eṉ moḻi}
  \item \emph{Recall:} say the Tamil for a family name, then the Tamil for a pet name, then say \emph{moḻi} again
\end{itemize}

\begin{tcolorbox}[breakable,colback=teal!4,colframe=teal!35!black,title={Before you move on}]
What does \ta{மொழி} mean? (\textquotedblleft{}A language\textquotedblright{}.) Say it once more.
\end{tcolorbox}
\section[\texorpdfstring{\ta{ஆங்கிலம்} (āṅkilam)}{ஆங்கிலம் (āṅkilam)}]{\ta{ஆங்கிலம்} (āṅkilam) — English, the language}

\label{lesson:TA-C104-english}



\begin{quote}
Before the new one: say the Tamil for a pet name, then the Tamil for a language.
\end{quote}

\subsection*{You'll want to know: \ta{ஆங்கிலம்}}
\textbf{\ta{ஆங்கிலம்}} (\emph{āṅkilam}) — \textquotedblleft{}English\textquotedblright{}.

\textbf{\ta{எனக்கு} \ta{ஆங்கிலம்} \ta{தெரியும்}} — \textquotedblleft{}I know English\textquotedblright{}.

\textbf{In the family, and next door}

\noindent
\begin{tabularx}{\linewidth}{@{}>{\raggedright\arraybackslash}X>{\raggedright\arraybackslash}X@{}}
\toprule
\textbf{} & \textbf{word} \\
\midrule
Kannada & \ka{ಇಂಗ್ಲಿಷ್} (\emph{iṅgliṣ}) \\
Telugu & \te{ఇంగ్లీషు} (\emph{iṅglīṣu}) \\
Malayalam & \ml{ഇംഗ്ലീഷ്} (\emph{iṁglīṣ}) \\
Hindi (neighbour) & \dv{अंग्रेज़ी} (\emph{angrezī}) \\
English & English, the language \\
\bottomrule
\end{tabularx}

\begin{grammarlens}[title={What you've built}]
English.
\end{grammarlens}

\subsection*{Your turn}
\begin{itemize}
\raggedright
  \item \emph{Say it:} \emph{āṅkilam}
  \item \emph{Say it:} \emph{āṅkilam}, once more
  \item \emph{Say it:} say \emph{eṉakku āṅkilam teriyum}
  \item \emph{Recall:} say the Tamil for a pet name, then the Tamil for a language, then say \emph{āṅkilam} again
\end{itemize}

\begin{tcolorbox}[breakable,colback=teal!4,colframe=teal!35!black,title={Before you move on}]
What does \ta{ஆங்கிலம்} mean? (\textquotedblleft{}English, the language\textquotedblright{}.) Say it once more.
\end{tcolorbox}
\section[\texorpdfstring{\ta{திருமணம்} (tirumaṇam)}{திருமணம் (tirumaṇam)}]{\ta{திருமணம்} (tirumaṇam) — a wedding, marriage}

\label{lesson:TA-C104-wedding}



\begin{quote}
Before the new one: say the Tamil for a language, then the Tamil for English.
\end{quote}

\subsection*{You'll want to know: \ta{திருமணம்}}
\textbf{\ta{திருமணம்}} (\emph{tirumaṇam}) — \textquotedblleft{}a wedding\textquotedblright{}, \textquotedblleft{}marriage\textquotedblright{}.

\textbf{\ta{திரு}-} is \textquotedblleft{}sacred, honoured\textquotedblright{}, as in the title \textbf{\ta{திரு}}, \textquotedblleft{}Mr\textquotedblright{}.

\textbf{In the family, and next door}

\noindent
\begin{tabularx}{\linewidth}{@{}>{\raggedright\arraybackslash}X>{\raggedright\arraybackslash}X@{}}
\toprule
\textbf{} & \textbf{word} \\
\midrule
Kannada & \ka{ಮದುವೆ} (\emph{maduve}) \\
Telugu & \te{పెళ్ళి} (\emph{peḷḷi}) \\
Malayalam & \ml{കല്യാണം} (\emph{kalyāṇaṁ}) \\
Hindi (neighbour) & \dv{शादी} (\emph{śādī}) \\
English & a wedding, marriage \\
\bottomrule
\end{tabularx}

\begin{grammarlens}[title={What you've built}]
A wedding.
\end{grammarlens}

\subsection*{Your turn}
\begin{itemize}
\raggedright
  \item \emph{Say it:} \emph{tirumaṇam}
  \item \emph{Say it:} \emph{tirumaṇam}, once more
  \item \emph{Say it:} say \emph{tirumaṇam}
  \item \emph{Recall:} say the Tamil for a language, then the Tamil for English, then say \emph{tirumaṇam} again
\end{itemize}

\begin{tcolorbox}[breakable,colback=teal!4,colframe=teal!35!black,title={Before you move on}]
What does \ta{திருமணம்} mean? (\textquotedblleft{}A wedding, marriage\textquotedblright{}.) Say it once more.
\end{tcolorbox}
\section[\texorpdfstring{\ta{முகவரி} (mukavari)}{முகவரி (mukavari)}]{\ta{முகவரி} (mukavari) — an address}

\label{lesson:TA-C104-address}



\begin{quote}
Before the new one: say the Tamil for English, then the Tamil for a wedding.
\end{quote}

\subsection*{You'll want to know: \ta{முகவரி}}
\textbf{\ta{முகவரி}} (\emph{mukavari}) — \textquotedblleft{}an address\textquotedblright{}, literally \textquotedblleft{}face-line\textquotedblright{}.

\textbf{\ta{உங்கள்} \ta{முகவரி} \ta{என்ன}?} — \textquotedblleft{}what is your address?\textquotedblright{}.

\textbf{In the family, and next door}

\noindent
\begin{tabularx}{\linewidth}{@{}>{\raggedright\arraybackslash}X>{\raggedright\arraybackslash}X@{}}
\toprule
\textbf{} & \textbf{word} \\
\midrule
Kannada & \ka{ವಿಳಾಸ} (\emph{viḷāsa}) \\
Telugu & \te{చిరునామా} (\emph{cirunāmā}) \\
Malayalam & \ml{വിലാസം} (\emph{vilāsaṁ}) \\
English & an address \\
\bottomrule
\end{tabularx}

\begin{grammarlens}[title={What you've built}]
An address.
\end{grammarlens}

\subsection*{Your turn}
\begin{itemize}
\raggedright
  \item \emph{Say it:} \emph{mukavari}
  \item \emph{Say it:} \emph{mukavari}, once more
  \item \emph{Say it:} ask \emph{uṅkaḷ mukavari eṉṉa?}
  \item \emph{Recall:} say the Tamil for English, then the Tamil for a wedding, then say \emph{mukavari} again
\end{itemize}

\begin{tcolorbox}[breakable,colback=teal!4,colframe=teal!35!black,title={Before you move on}]
What does \ta{முகவரி} mean? (\textquotedblleft{}An address\textquotedblright{}.) Say it once more.
\end{tcolorbox}
\section[\texorpdfstring{\ta{பிறந்தநாள்} (piṟantanāḷ)}{பிறந்தநாள் (piṟantanāḷ)}]{\ta{பிறந்தநாள்} (piṟantanāḷ) — a birthday}

\label{lesson:TA-C104-birthday}



\begin{quote}
Before the new one: say the Tamil for a wedding, then the Tamil for an address.
\end{quote}

\subsection*{You'll want to know: \ta{பிறந்தநாள்}}
\textbf{\ta{பிறந்தநாள்}} (\emph{piṟantanāḷ}) — \textquotedblleft{}a birthday\textquotedblright{}, literally \textquotedblleft{}the day one was born\textquotedblright{}.

\textbf{\ta{பிறந்தநாள்} \ta{வாழ்த்துகள்}} — \textquotedblleft{}happy birthday\textquotedblright{}, with the wishes you know.

\textbf{In the family, and next door}

\noindent
\begin{tabularx}{\linewidth}{@{}>{\raggedright\arraybackslash}X>{\raggedright\arraybackslash}X@{}}
\toprule
\textbf{} & \textbf{word} \\
\midrule
Kannada & \ka{ಹುಟ್ಟುಹಬ್ಬ} (\emph{huṭṭuhabba}) \\
Telugu & \te{పుట్టినరోజు} (\emph{puṭṭinarōju}) \\
Hindi (neighbour) & \dv{जन्मदिन} (\emph{janmdin}) \\
English & a birthday \\
\bottomrule
\end{tabularx}

\begin{grammarlens}[title={What you've built}]
Language, English, wedding, address, birthday.
\end{grammarlens}

\subsection*{Your turn}
\begin{itemize}
\raggedright
  \item \emph{Say it:} \emph{piṟantanāḷ}
  \item \emph{Say it:} \emph{piṟantanāḷ}, once more
  \item \emph{Say it:} say \emph{piṟantanāḷ vāḻttukaḷ}
  \item \emph{Recall:} say the Tamil for a wedding, then the Tamil for an address, then say \emph{piṟantanāḷ} again
\end{itemize}

\begin{tcolorbox}[breakable,colback=teal!4,colframe=teal!35!black,title={Before you move on}]
What does \ta{பிறந்தநாள்} mean? (\textquotedblleft{}A birthday\textquotedblright{}.) Say it once more.
\end{tcolorbox}
